\subsection{LOWER CENTRAL SERIES}

PROPOSITION 5. Let \(\mathfrak{g}\) be a Lie algebra and \(\mathrm{P}\) a submodule of \(\mathfrak{g}\) . We define the submodules \(\mathrm{P}_n\) of \(\mathfrak{g}\) by the formulae \(\mathrm{P}_1 = \mathrm{P}\) and \(\mathrm{P}_{n+1} = [\mathrm{P},\mathrm{P}_n]\) for \(n\geqslant 1\) . Then

\begin{enumerate}
\def\labelenumi{(\arabic{enumi})}
\setcounter{enumi}{17}
\tightlist
\item
\end{enumerate}

\[
[ \mathrm{P} _ {m}, \mathrm{P} _ {n} ] \subset \mathrm{P} _ {m + n},\tag{18}
\]

\[
\mathrm{P} _ {n} = \sum_ {p = 1} ^ {n - 1} \left[ \mathrm{P} _ {p}, \mathrm{P} _ {n - p} \right] for n \geqslant 2.\tag{19}
\]

We prove (18) by induction on \(m\) . The case \(m = 1\) is obvious. By the Jacobi identity,

\[
[ [ \mathrm{P}, \mathrm{P} _ {m} ], \mathrm{P} _ {n} ] \subset [ \mathrm{P} _ {m}, [ \mathrm{P}, \mathrm{P} _ {n} ] ] + [ \mathrm{P}, [ \mathrm{P} _ {m}, \mathrm{P} _ {n} ] ],
\]

that is

\[
[ \mathrm{P} _ {m + 1}, \mathrm{P} _ {n} ] \subset [ \mathrm{P} _ {m}, \mathrm{P} _ {n + 1} ] + [ \mathrm{P}, [ \mathrm{P} _ {m}, \mathrm{P} _ {n} ] ].
\]

The induction hypothesis implies that \([\mathbf{P}_m, \mathbf{P}_{n+1}] \subset \mathbf{P}_{m+n+1}\) and \([\mathbf{P}_m, \mathbf{P}_n] \subset \mathbf{P}_{m+n}\) , whence

\[
[ \mathrm{P} _ {m + 1}, \mathrm{P} _ {n} ] \subset \mathrm{P} _ {m + n + 1} + [ \mathrm{P}, \mathrm{P} _ {m + n} ] = \mathrm{P} _ {m + n + 1}.
\]

By formula (18), \(\mathrm{P}_n\supset \sum_{p = 1}^{n - 1}[\mathrm{P}_p,\mathrm{P}_{n - p}]\supset [\mathrm{P}_1,\mathrm{P}_{n - 1}] = \mathrm{P}_n\) , whence (19).

When we take \(\mathbf{P} = \mathfrak{g}\) , the sequence \((\mathrm{P}_n)\) is the lower central series \((\mathcal{C}^n\mathfrak{g})\) of \(\mathfrak{g}\) (Chapter I, § 1, no. 5). Hence:

PROPOSITION 6. Let g be a Lie algebra and \((\mathcal{C}^{n}\mathfrak{g})_{n\geqslant 1}\) the lower central series of g. Then

\[
[ \mathcal {C} ^ {m} \mathfrak {g}, \mathcal {C} ^ {n} \mathfrak {g} ] \subset \mathcal {C} ^ {m + n} \mathfrak {g} \quad for m \geqslant 1 and n \geqslant 1.
\]

Generalizing Definition 1 of Chapter I, §4, no. 1, we shall say that a Lie algebra g is nilpotent if \(\mathcal{C}^{n}\mathfrak{g} = \{0\}\) for \(n\) sufficiently large. The nilpotency class of a nilpotent Lie algebra g is the smallest integer \(n\) such that \(\mathcal{C}^{n + 1}\mathfrak{g} = \{0\}\) .

PROPOSITION 7. Let X be a set and n an integer ≥1.

\begin{enumerate}
\def\labelenumi{(\alph{enumi})}
\item
  \(\mathrm{L}^{n+1}(\mathrm{X}) = [\mathrm{L}^{1}(\mathrm{X}), \mathrm{L}^{n}(\mathrm{X})]\) .
\item
  The module \(L^{n}(X)\) is generated by the elements \([x_{1}, [x_{2}, \ldots, [x_{n-1}, x_{n}] \ldots]]\) where \((x_{1}, \ldots, x_{n})\) runs through the set of sequences of n elements of X.
\item
  The lower central series of L(X) is given by \(\mathcal{C}^n (\mathrm{L}(\mathrm{X})) = \sum_{p\geqslant n}\mathrm{L}^p (\mathrm{X})\)
\item
  We apply Proposition 5 with \(\mathfrak{g} = \mathrm{L}(\mathrm{X})\) and \(\mathrm{P} = \mathrm{L}^{1}(\mathrm{X})\) . By induction on \(n\) , we deduce from (12) (no. 6) and (19) the equality \(\mathrm{P}_{n} = \mathrm{L}^{n}(\mathrm{X})\) . The desired relation is then equivalent to the definition \([\mathrm{P}, \mathrm{P}_{n}] = \mathrm{P}_{n + 1}\) .
\item
  This follows from (a) by induction on \(n\) .
\item
  Let \(\mathfrak{g} = \mathrm{L}(\mathrm{X})\) and \(\mathfrak{g}_{n} = \sum_{p\geqslant n}\mathrm{L}^{p}(\mathrm{X})\) . Then \(\mathfrak{g} = \mathfrak{g}_1\) and formula (13) of no. 6 implies \([\mathfrak{g}_{n},\mathfrak{g}_{m}]\subset \mathfrak{g}_{n + m}\) and in particular \([\mathfrak{g},\mathfrak{g}_n]\subset \mathfrak{g}_{n + 1}\) . By induction on \(n,\mathcal{C}^{n}\mathfrak{g}\subset \mathfrak{g}_{n}\) . On the other hand, from (a) we deduce that \(\mathrm{L}^{n}(\mathrm{X})\subset \mathcal{C}^{n}\mathfrak{g}\) by induction on \(n\) . As \(\mathcal{C}^{n}\mathfrak{g}\) is an ideal of \(\mathfrak{g}\) , the relation \(\mathrm{L}^p (\mathrm{X})\subset \mathcal{C}^{n}\mathfrak{g}\) implies that
\end{enumerate}

\[
\mathrm{L} ^ {p + 1} (\mathrm{X}) = [ \mathrm{L} ^ {1} (\mathrm{X}), \mathrm{L} ^ {p} (\mathrm{X}) ] \subset \mathcal {C} ^ {n} \mathfrak {g}
\]

by (a). Hence \(\mathbf{L}^p (\mathbf{X})\subset \mathcal{C}^{n}\mathfrak{g}\) for \(p\geqslant n\) , whence \(\mathfrak{g}_n\subset \mathcal{C}^{n}\mathfrak{g}\)

COROLLARY. Let \(\mathfrak{g}\) be a Lie algebra and \((x_{i})_{i\in I}\) a generating family of \(\mathfrak{g}\) . The \(n\) -th term \(\mathcal{C}^{n}\mathfrak{g}\) of the lower central series of \(\mathfrak{g}\) is the module generated by the iterated brackets \([x_{i_1}, [x_{i_2}, \ldots, [x_{i_{p-1}}, x_{i_p}] \ldots]]\) for \(p \geqslant n\) , and \(i_1, \ldots, i_p\) in I.

Let \(f\) be the homomorphism of L(I) into g such that \(f(i) = x_i\) for all \(i \in I\) . As \((x_i)_{i \in I}\) generates g, g = f(L(I)), whence \(\mathcal{C}^{\bullet}g = f(\mathcal{C}^{\bullet}(L(I)))\) by Proposition 4 of Chapter I, §1, no. 5. The corollary then follows from assertions (b) and (c) of Proposition 7.

